
\SourceLeaf{492}\Article{29}
Als ein besonders merkwürdiges und willkommenes Resultat erscheint die Kleinheit der für $e,e',e'',e'''$, oder vielmehr zunächst für ihre Unterschiede von ihrem Mittel $\varepsilon$ gefundenen Werthe. Eine besondere Untersuchung hat das Gewicht dieser Bestimmungen $\frac{96}{11}$ mal grösser als das Gewicht von $\tfrac12(f+g)$ ergeben, folglich die mittlere daran haftende Unsicherheit $=60''5\BareRoot{\frac{11}{96}}=20''5$, woraus erhellt, dass sogar die Realität von Ungleichheiten zwischen $e,e',e'',e'''$ ganz zweifelhaft bleibt. Da es nun höchst unwahrscheinlich ist, dass bei vier Nadeln constante Fehler von fast genau gleicher Grösse Statt finden sollten, so ist man berechtigt anzunehmen, dass dieselben gar keine oder doch nur ganz unmerkliche constante Fehler haben, und es möchte daher fast unnöthig scheinen, von der Drehbarkeit der Achsen an zweien derselben zu weitern Proben einen Gebrauch zu machen.

Für eine der Nadeln, nemlich für Nr.\ 4, geben wirklich schon einige frühere Beobachtungen eine Verstärkung dieses Schlusses. Es waren nemlich an vier Tagen vom 15.—19. Mai mit den Nadeln 3 und 4 ähnlich combinirte Beobachtungen gemacht, wie später vom 8.—25. Junius, nur mit dem Unterschiede, dass jedes partielle Resultat nicht auf acht, sondern nur auf vier Einstellungen beruhte; an der Nadel 3 waren die Zapfen in derselben Lage wie später, aber an der Nadel 4 standen sie anders, indem nach dem 19. Mai eine Drehung von etwa einem Quadranten vorgenommen ist. Die Beobachtungen, eben so geschrieben wie im 23. Art., sind folgende:

\Needspace{8\baselineskip}
\begin{center}\textit{Beobachtungen mit Nadel} 3.

\begin{tabular}{rr|c|rrr|rrr|l}
Mai&15&$B$&$67^0$&$41'$&$26''$&$67^0$&$44'$&$53''$&$6''16166$\\
&17&$B$&&$43$&$52$&&$45$&$52$&$6.20333$\\
&18&$A$&&$33$&$56$&&$39$&$15$&$6.17781$\\
&19&$A$&&$36$&$8$&&$37$&$8$&$6.19566$
\end{tabular}
\end{center}
\Needspace{8\baselineskip}
\begin{center}\textit{Beobachtungen mit Nadel} 4.

\begin{tabular}{rr|c|rrr|rrr|l}
Mai&15&$A$&$67^0$&$14'$&$28''$&$67^0$&$47'$&$49''$&$5''94332$\\
&17&$B$&$68$&$5$&$39$&&$36$&$36$&$5.92034$\\
&18&$B$&&$3$&$30$&&$36$&$13$&$5.94235$\\
&19&$A$&$67$&$3$&$4$&&$59$&$47$&$5.94663$
\end{tabular}
\end{center}

Die Beobachtungen sind alle in den Vormittagsstunden gemacht.

\SourceLeaf{493}
Zur Berechnung sind bei Nadel $3$ die oben gefundenen Werthe von $x'',y'',e''$ angewandt; bei Nadel $4$ mussten hingegen die Werthe von $x''',y''',e'''$, so gut es angeht, aus diesen Beobachtungen selbst abgeleitet werden, wobei gefunden wurde
\[
\begin{aligned}
y'''&=-1103''\\
x'''&=+556''\\
e'''-\varepsilon&=+24''
\end{aligned}
\]
Die Bestimmung von $y'''$, auf so wenige Beobachtungen gegründet, ist allerdings sehr unsicher, allein der Einfluss davon auf die Reduction von $\tfrac12(f+g)$ bleibt ganz unbedeutend, indem der grösste Coëfficient von $y'''$ in den Gleichungen I nur $0.00341$ ist. Die Resultate für $i$ stehen dann so:
\begin{center}
\begin{tabular}{rr||rrr||rrr}
\multicolumn{2}{c||}{}&\multicolumn{3}{c||}{Nadel $3$}&\multicolumn{3}{c}{Nadel $4$}\\\hline
Mai&15&$67^0$&$38'$&$57''$&$67^0$&$40'$&$0''$\\
&17&&$40$&$36$&&$41$&$36$\\
&18&&$41$&$15$&&$40$&$15$\\
&19&&$41$&$19$&&$40$&$16$
\end{tabular}
\end{center}
Das Gewicht der Bestimmung von $e'''-\varepsilon$ wird hier nur doppelt so gross, als das Gewicht von $\tfrac12(f+g)$, und da die Beobachtungen selbst eine bedeutend geringere Genauigkeit haben, als die spätern, so erhellt, dass der jetzt gefundene Werth eben so wenig für die Realität eines constanten Fehlers spricht, als der aus den spätern Beobachtungen abgeleitete.

Die starke Abweichung der Werthe von $x'''$ und $y'''$ von den oben (Art.\ 26.\ 27) gefundenen, beweist nur, dass der drehbare Theil der Nadel für sich betrachtet seinen Schwerpunkt nicht in der Zapfenachse hat, woran übrigens auch wenig gelegen ist.

\Article{30}
Ich stelle nun noch die Endresultate für die Inclination aus den sämmtli\-chen behandelten Beobachtungen zusammen, und nehme unter dieselben auch die Resultate der schon oben erwähnten Beobachtungen vom 1. und 7. August mit auf, welche mit der Nadel 1 ganz auf dieselbe Art wie am 23. September gemacht sind. Diese Beobachtungen selbst waren folgende:
\SourceLeaf{494}
\begin{center}
\begin{tabular}{r||rrr||rrr}
&\multicolumn{3}{c||}{August $1$}&\multicolumn{3}{c}{August $7$}\\\hline
$f$&$67^0$&$20'$&$12''$&$67^0$&$22'$&$41''$\\
$g$&&$44$&$11$&&$42$&$8$\\
$f'$&&$59$&$53$&$68$&$1$&$56$\\
$g'$&&$35$&$53$&$67$&$35$&$46$
\end{tabular}
\end{center}
\Needspace{20\baselineskip}
\begin{center}\textit{Inclinationsbestimmungen}

\begin{tabular}{rrr|rrr||rr|rrr}
1842&Mai&15&$67^0$&$39'$&$28''$&Juni&18&$67^0$&$41'$&$0''$\\
&&17&&$41$&$6$&&22&&$37$&$58$\\
&&18&&$40$&$45$&&23&&$40$&$4$\\
&&19&&$40$&$47$&&25&&$39$&$47$\\
&&20&&$40$&$18$&Juli&6&&$41$&$0$\\
&&21&&$39$&$26$&&7&&$40$&$30$\\
&&22&&$39$&$36$&&8&&$40$&$4$\\
&&24&&$39$&$2$&&9&&$39$&$34$\\
&&31&&$39$&$47$&&17&&$40$&$1$\\
&Juni&2&&$37$&$27$&&18&&$39$&$44$\\
&&4&&$37$&$15$&&19&&$39$&$4$\\
&&5&&$40$&$52$&&20&&$39$&$22$\\
&&8&&$43$&$43$&Aug.&1&&$39$&$57$\\
&&9&&$37$&$51$&&7&&$40$&$26$\\
&&11&&$37$&$14$&Sept.&23&&$40$&$54$\\
&&16&&$37$&$42$&&&&&
\end{tabular}
\end{center}
Das Mittel aus allen $31$ Bestimmungen, ohne einen Gewichtsunterschied zu berücksichtigen, wird
\[
67^0\;39'\;44''
\]
und mag als für den 21. Junius gültig angesehen werden. Das Mittel aus den $24$ Bestimmungen vom 20. Mai bis 20. Julius allein, dem als mittlerer Zeitpunkt der 19. Junius entspricht, ist
\[
67^0\;39'\;31''
\]
\Article{31}
Die Unterschiede der Inclinationen für die einzelnen $31$ Tage von ihrem Mittel sind zusammengesetzt aus der noch nachbleibenden Wirkung der Beobachtungsfehler und den wirklichen Ungleichheiten der Inclination selbst. Für die
\SourceLeaf{495}einzelnen Tage lassen sich zwar diese Bestandtheile nicht von einander scheiden, allein eine Abschätzung eines Mittelwerths der wirklichen Schwankungen mag bei einer so zahlreichen Reihe wohl versucht werden. In dieser Absicht habe ich zuvörderst die Inclinationen unter Voraussetzung einer regelmässigen jährlichen Abnahme von $3$ Minuten auf den 21. Junius reducirt, und dann die Quadrate der Differenzen von dem Mittelwerthe addirt; diese Summe $220184$ mit $30$ dividirt gibt $7339.5$ als Quadrat des mittlern Fehlers, dem man sich aussetzt, wenn man aufs Gerathewohl eine jener $31$ Inclinationen als die mittlere für die Zeit der Beobachtung gültige ansehen wollte. Soll die ungleiche Zuverlässigkeit der drei Beobachtungsgruppen berücksichtigt werden, so ergeben die Grundsätze der Wahrscheinlichkeitsrechnung, indem man den mittlern Fehler für die vier ersten Beobachtungen mit $m'$, für die drei letzten mit $m''$, und für die $24$ übrigen mit $m$, das mittlere Schwanken der Inclination selbst aber mit $M$ bezeichnet, folgende Gleichung:
\[
7339{,}5=\frac{24mm+4m'm'+3m''m''}{31}+MM
\]
Für $mm$ ist oben der Werth $1829.25$ gefunden, oder es kann wenigstens diese Zahl wie eine hinlängliche Annäherung angesehen werden, für die sieben andern Beobachtungen mag in Ermangelung eines sichern Maassstabes die Zahl der Einstellungen, woraus die Resultate abgeleitet sind, zum Grunde gelegt, also
\[
m'm'=2mm,\qquad m''m''=\tfrac85mm
\]
gesetzt werden. Dadurch wird
\[
MM=7339{,}5-\tfrac{184}{155}.1829{,}25=5168
\]
und $M=71''9$.

\Article{32}
Mit demselben Instrumente und an demselben Platze hatte ich auch schon im vorigen Jahre eine Reihe von Inclinationsbeobachtungen gemacht, von denen ich jedoch nur die Endresultate hieher setze.
\SourceLeaf{496}
\begin{center}
\begin{tabular}{rrrcrrr}
1841&Sept.&22&$.\quad.\quad.$&$67^0$&$40'$&$20''$\\
&&24&&&$40$&$53$\\
&&27&&&$46$&$41$\\
&Oct.&2&&&$42$&$57$\\
&&7&&&$42$&$14$\\
&&10&&&$42$&$40$\\
&&12&&&$43$&$15$\\
&&20&&&$44$&$2$\\
&&20&&&$42$&$5$\\
&&22&&&$42$&$52$\\\hline
Mittel,&Oct.&8&&$67^0$&$42'$&$48''$
\end{tabular}
\end{center}
Die ersten acht Beobachtungen sind auf ähnliche Art angestellt, wie die diesjährigen, indem an jedem Tage, ohne die Pole zwischen den Beobachtungen umzukehren, zwei Nadeln (Nr.\ 1 und 2) angewandt wurden; die beiden letzten hingegen wurden auf die gewöhnliche Art gemacht, die zweite vom 20. Oct. mit Nadel 4, die vom 22. mit Nadel 3. Die Zeit war am 27. Sept. und 10. Oct. Nachmittags zwischen $3$ und $5$ Uhr, bei allen übrigen Vormittags. Jede dieser $10$ Inclinationen beruhete auf $16$ Einstellungen, und es wird ihnen aus diesem Grunde auch nur ein verhältnissmässig kleineres Gewicht zuzuerkennen sein, als den Inclinationen von $1842$, die resp.\ auf $32$, $64$ und $40$ Einstellungen beruheten.

\Article{33}
Sämmtliche bisher angeführte Inclinationen bedürfen noch einer kleinen gemeinschaftlichen Correction wegen des Einflusses, welchen an dem Beobachtungsplatze die Magnetstäbe der Magnetometer in der Sternwarte und im magnetischen Observatorium ausüben. Um die Resultate davon zu befreien, muss durchgehens $5''15$ abgezogen werden (vergl.\ \textit{Resultate} für 1840. II. Art.\ 6) [S.\ 433 d.\ B.]

Die \textit{absolute} Zuverlässigkeit der Inclinationsbestimmungen bleibt übrigens noch abhängig von der Richtigkeit der Voraussetzung, dass das Instrument selbst keine Theile enthält, die eine magnetische Wirkung auf die Nadel haben können. Ein Grund zu einer solchen Befürchtung ist bei dem von mir gebrauchten Instrumente nicht vorhanden; einige Beobachtungen, die ich nach der im 18. Art. erwähnten Art mit einer belasteten Nadel anstellte, haben immer nur Abweichun\SourceLeaf{497}gen von ein Paar Minuten gezeigt, die sich aus den unvermeidlichen zufälligen Beobachtungsfehlern und den wirklichen Anomalien der Inclination selbst ganz ungezwungen erklären lassen. Auch die hinlänglich befriedigende Übereinstimmung der Werthe, welche im 11. Art. für die daselbst mit $\alpha$ bezeichnete Grösse gefunden sind, spricht gegen das Vorhandensein von solchen Störungen. Zur Erkennung ganz kleiner Einflüsse sind freilich solche Prüfungen nicht geeignet, und ich muss mir daher die weitere Prüfung durch mehr durchgreifende Mittel vorbehalten.

\Article{34}
Zum Schluss stelle ich noch meine Resultate mit einigen ältern Bestimmungen zusammen.
\[
\begin{array}{r|l}
\begin{array}{rrr}1805&\text{Dec.}&\\1826&\text{Sept.}&\end{array}&
\left.\begin{array}{rrr}69^0&29'&\\68&29&26''\end{array}\right\}\ \text{\textsc{von Humboldt}}\\
\begin{array}{rrr}1837&\text{Juli}&1\\&\text{—}&\text{—}\end{array}&
\left.\begin{array}{rrr}67&47&0\\67&53&30\end{array}\right\}\ \text{\textsc{Forbes}}\\
\begin{array}{rrr}1841&\text{Oct.}&8\\1842&\text{Juni}&21\end{array}&
\begin{array}{rrr}67&42&43\\67&39&39\end{array}
\end{array}
\]
Die beiden ersten Beobachtungen habe ich aus den \textit{Additions} zu dem XIII. Bande der \textit{Voyage aux régions équinoxiales} entlehnt (S.\ 152); die erste ist mit einem Inclinatorium von \textsc{Lenoir}, die zweite mit einem Instrument von \textsc{Gambey} angestellt; letztere beruhet auf den Beobachtungen mit zwei Nadeln, deren Resultate a.\ a.\ O. zu $68^0\;30'\;7''$ und $68^0\;28'\;15''$ angegeben werden, womit das ebendaselbst angesetzte Mittel nicht übereinstimmt; vermuthlich ist die Zahl für die zweite Nadel durch einen Druckfehler um $30''$ zu klein angesetzt. Der Beobachtungsplatz $1805$ ist mir nicht bekannt; $1826$ war er im freien Felde einige hundert Schritte östlich von der Sternwarte.

\textsc{Forbes} Beobachtungen sind in den \textit{Transactions of the Royal Society of Edinburgh} Vol.\ XV, Part.\ 1, S.\ 31 und 32 abgedruckt; sie wurden an einem \textsc{Robinson}schen Instrument von kleinern Dimensionen als das hiesige mit zwei Nadeln von $6$ engl.\ Zoll Länge im Garten der Sternwarte angestellt; die zweite Nadel hält der Beobachter selbst für die bessere.

Ich habe unter diese Beobachtungen die von \textsc{Mayer} im März 1814 angestellten und in den \textit{Commentationes recent.\ Soc.\ Gotting.} T.\ III, S.\ 36 u.\ 37 ange\SourceLeaf{498}führten nicht einreihen wollen, da dieselben gar kein Vertrauen verdienen. Wie sehr unvollkommen das von \textsc{Mayer} gebrauchte Instrument war, zeigt die von ihm selbst S.\ 35 gegebene Probe, wo bei bleibender Stellung des Instruments zehn wiederholte Einstellungen Differenzen von mehr als einem Grade gaben. Seine Resultate für die Inclination selbst, von zwei verschiedenen Tagen, weichen um einen halben Grad von einander ab.

Eben so wenig verdiente meine eigne Beobachtung vom 23. Juni 1832, die in der \textit{Intens.\ vis magneticae terrestris} art.\ 27 angeführt ist, hier einen Platz, sowohl wegen der Unvollkommenheit des Instruments, als wegen des Locals in der Sternwarte, wo nicht sehr entferntes Eisenwerk das Resultat bedeutend afficiren, und zwar nachweislich eine Vergrösserung der Inclination hervorbringen musste.

Die angeführten Inclinationen lassen sich nun zwar sehr gut durch die Annahme einer jährlichen gleichförmigen Verminderung von $3$ Minuten oder genauer $3'\;2''3$ vereinigen, wenn man bei \textsc{Forbes} Beobachtungen sich an das Resultat der zweiten Nadel hält, und es bleiben nur Abweichungen übrig, die füglich dem Conspiriren der Beobachtungsfehler und der Schwankungen der Inclination zugeschrieben werden können. Da jedoch nach \textsc{Hansteens} Untersuchungen über die Beobachtungen an andern europäischen Orten die jährliche Abnahme allmählig langsamer geworden ist, so wird man die angegebene Zahl nur wie einen mittlern etwa für $1829$ gültigen Werth zu betrachten, und die Bestätigung und genauere Festsetzung der Ungleichförmigkeit erst von künftigen Beobachtungen zu erwarten haben.

\clearpage
\SourceLeaf{499}\pdfbookmark[0]{Aufsätze über verschiedene Gegenstände der mathematischen Physik.}{various-physics-title}
\thispagestyle{empty}
\vspace*{0.23\textheight}
\begin{center}
{\Large AUFSÄTZE\par}
\vspace{2\baselineskip}
{\large ÜBER VERSCHIEDENE GEGENSTÄNDE\par}
\vspace{2\baselineskip}
{\Large\bfseries DER MATHEMATISCHEN PHYSIK.\par}
\end{center}
\clearpage
\fancyhead[L]{\small FUNDAMENTALGLEICHUNGEN.}
\SourceLeaf{501}\pdfbookmark[0]{Fundamentalgleichungen für die Bewegung schwerer Körper auf der rotirenden Erde.}{rotating-earth-title}
\begin{center}
{\large\bfseries FUNDAMENTALGLEICHUNGEN\par\medskip FÜR DIE BEWEGUNG SCHWERER KÖRPER\par\medskip AUF DER ROTIRENDEN ERDE.\par}
\bigskip
\textsc{Benzenberg}. Versuche über das Gesetz des Falls. 1804.
\par\bigskip
{\large\itshape Brief von Gauss an Benzenberg.\par}
\end{center}
\begin{flushright}Braunschweig 1803. Februar 2.\end{flushright}
—\quad—\quad—\quad In der Theorie unsres Freundes \textsc{Olbers} ist eine Voraussetzung, die mir nicht zulässig scheint. Nemlich: \textit{dass der Körper während des Falls in einer Ebene bleibe.} Allein dies darf man, meiner Meinung nach, \textit{nicht} voraussetzen, wenn man den Widerstand der Luft in Betracht zieht, den man hier \textit{nothwendig} in Betracht ziehen muss, weil die geschlossene Abweichung nach Süden lediglich darauf beruht. Eine leichte Betrachtung zeigt nemlich folgendes: die Ebene $(\mathrm{A})$, in welcher der Körper sich ursprünglich zu bewegen anfängt, geht durch den Mittelpunkt der Erde (oder allgemeiner, der Attraction), und steht auf derjenigen Ebene $(\mathrm{B})$ senkrecht, in der der Meridian des Beobachtungsorts beim Anfang des Falls war. Allein man sieht leicht, dass die Lufttheile an allen Stellen der Ebene $\mathrm{A}$ schief dadurch gehen, bloss die gerade Linie ausgenommen, wo $\mathrm{A}$ von $\mathrm{B}$ geschnitten wird. Die Luft wirkt daher dem Körper nicht in dieser Ebene $\mathrm{A}$ entgegen, sondern treibt ihn daraus weg nach Norden, und es schien mir, dass der Effect davon gerade so gross sein würde, dass er die aus der Verspätung des Falls geschlossene Abweichung nach Süden aufhöbe.

Nachdem ich durch Ihren letzten Brief veranlasst war, aufs Neue an diese Materie zu denken, betrachtete ich in einer müssigen halben Stunde die Sache auf eine ganz verschiedene Art, und entwickelte die analytischen Gleichungen,
